\exercisesfor{9}

\begin{enumerate}
\def\labelenumi{\arabic{enumi}.}
\item
  Let \(\mathrm{G} = \mathbf{SO}(3, \mathbf{R})\) . Let \(\Delta_{1}, \Delta_{2}\) be two orthogonal lines in \(R^{3}\) and \(H_{i}\) the subgroup of G consisting of the rotations about \(\Delta_{i} (i = 1, 2)\) . Then \([\mathrm{L}(\mathrm{H}_{1}), \mathrm{L}(\mathrm{H}_{2})]\) is a 1-dimensional Lie subalgebra of G distinct from the Lie subalgebra tangent to \((\mathrm{H}_{1}, \mathrm{H}_{2})\) at e.
\item
  Suppose that K is ultrametric. Let G be a finite-dimensional Lie group, g its Lie algebra and A a finite subset of G. Then \(Z_{G}(A)\) is a Lie subgroup of G with Lie algebra \(\mathfrak{z}_{\mathfrak{g}}(A)\) . (Argue as in Proposition 8.)
\item
  Let G be a connected real or complex Lie group. The centre Z of G is a Lie quasi-subgroup of G and L(Z) is the centre of L(G).
\item
  Suppose that K is ultrametric and p \textgreater{} 0 (in the notation of §7). Let G be a finite-dimensional Lie group, A a group of automorphisms of G and B the corresponding group of automorphisms of L(G). Let \(\mathrm{G}^{\mathbf{A}}\) (resp. \(\mathrm{L}(\mathrm{G})^{\mathbf{B}}\) ) be the set of elements of G (resp. L(G)) which are fixed under A (resp. B). Then \(G^{A}\) is a Lie subgroup of G with Lie algebra \(\mathrm{L}(G)^{\mathbf{B}}\) . (Use the logarithmic mapping.)
\item
  Let \(G\) be a finite-dimensional connected real or complex Lie group. Let \((G_0, G_1, \ldots)\) be the upper central series of (Chapter II, § 4, Exercise 18) and let \((g_0, g_1, \ldots)\) be the upper central series of the Lie algebra \(L(G)\) (Chapter I, § 1, no. 6). Then, for all \(i\) , \(G_i\) is a Lie subgroup of \(G\) such that \(L(G_i) = g_i\) .
\item
  Let \(\Gamma\) be the 3-dimensional simply connected nilpotent real Lie group defined in Exercise 5 (b) of §4. Let \(\alpha \in \mathbf{R}\) be an irrational number. Let \(P\) be the discrete subgroup of \(\Gamma \times \mathbf{R}\) consisting of the \(((0,0,x),\alpha x)\) , where \(x \in \mathbf{Z}\) . Let \(G = (\Gamma \times \mathbf{R}) / P\) . Show that \((G,G)\) is not closed in \(G\) .
\item
  \begin{enumerate}
  \def\labelenumii{(\alph{enumii})}
  \tightlist
  \item
    Let \(G\) be a semi-simple connected real Lie group, \(Z\) its centre and \(\rho\) a continuous linear representation of \(G\) on a finite-dimensional complex vector space. There exists an integer \(p\) such that, for all \(z \in Z\) , \(\rho(z)\) is diagonalizable and all the eigenvalues of \(\rho(z)\) are \(p\) -th roots of unity. (It can be assumed that \(\rho\) is irreducible. Then \(\rho(z)\) is scalar by Schur's Lemma. On the other hand, det \(\rho(g) = 1\) for all \(g \in G\) because \(G = \mathcal{D}G\) .)
  \end{enumerate}
\end{enumerate}

\begin{enumerate}
\def\labelenumi{(\alph{enumi})}
\setcounter{enumi}{1}
\tightlist
\item
  Deduce that, if \(G\) admits an injective finite-dimensional continuous linear representation, \(Z\) is finite.
\end{enumerate}

¶ 8. Let G be a finite-dimensional connected real Lie group, g its Lie algebra, n the largest nilpotent ideal of g and h = {[}g, g{]} + n.

\begin{enumerate}
\def\labelenumi{(\alph{enumi})}
\item
  \(\mathfrak{h}\) is a characteristic ideal of \(\mathfrak{g}\) ; the radical of \(\mathfrak{h}\) is \(n\) ; for all \(x \in \mathfrak{h}\) , \(\operatorname{Tr} \operatorname{ad}_{\mathfrak{h}} x = 0\) ; for every Levi section \(l\) of \(\mathfrak{g}\) , \(\mathfrak{h} = l + n\) .
\item
  Suppose that G is simply connected. Let Z be its centre, H the Lie subgroup of G with Lie algebra b and \(\phi\) the canonical morphism of G onto G/H. Then \(\phi(Z)\) is discrete in G/H. (The group G is the semi-direct product of a semi-simple Lie group S and its radical R. Let \(z \in Z\) . Then \(z = y^{-1}x\) , where \(x \in R\) and y belongs to the centre of S. There exists an integer p such that the eigenvalues of \(Ad_{q}y\) , and hence also of \(Ad_{q}x\) , are p-th roots of unity (Exercise 7). Deduce that, if N is the Lie subgroup of G with Lie algebra n, there exists a neighbourhood U of e in R satisfying U = UN = NU,
\end{enumerate}

\[
\mathrm{Z} \cap (\mathrm{SU}) \subset \mathrm{SN} = \mathrm{H}.
\]

\begin{enumerate}
\def\labelenumi{(\alph{enumi})}
\setcounter{enumi}{2}
\item
  We no longer assume that G is simply connected. Let H be the integral subgroup of G with Lie algebra b. Show that H is a unimodular normal Lie subgroup of G with nilpotent radical, such that G/H is commutative. (Use (a) and (b).)
\item
  Deduce that, if (G, G) is dense in G, the radical of G is nilpotent.
\end{enumerate}

\begin{enumerate}
\def\labelenumi{\arabic{enumi}.}
\setcounter{enumi}{8}
\tightlist
\item
  Let \(G\) be the universal covering of \(\mathbf{SL}(2, \mathbf{R})\) . We identify the kernel of the canonical morphism \(G \to \mathbf{SL}(2, \mathbf{R})\) with \(\mathbf{Z}\) (§ 6, Exercise 2). Let \(a\) be an element of \(\mathbf{T}^n\) whose powers are everywhere dense in \(\mathbf{T}^n\) (General Topology, Chapter VII, §1, Corollary 2 to Proposition 7). Let \(D\) be the discrete subgroup of \(G \times \mathbf{T}^n\) generated by (1, a). Let \(H = (G \times \mathbf{T}^n)/D\) . Then
\end{enumerate}

\[
\mathrm{L} (\mathrm{H}) = \mathfrak {s l} (2, \mathbf {R}) \times \mathbf {R} ^ {n}
\]

and the integral subgroup \(\mathbf{H}'\) of \(\mathbf{H}\) with Lie algebra \(\mathfrak{sl}(2,\mathbf{R})\) is isomorphic to \(\mathbf{G}\) and dense in \(\mathbf{H}\) . \(\mathrm{D}^n\mathrm{H}' = \mathrm{H}'\) for all \(n\geqslant 0\) .

¶ 10. Let G be a finite-dimensional real or complex Lie group. Let

\[
p = \dim [ L (G), L (G) ].
\]

Let \((\mathrm{G},\mathrm{G})\) be given its structure as an integral subgroup of \(\mathbf{G}\) . There exists a neighbourhood \(\mathbf{V}\) of \(e\) in \((\mathrm{G},\mathrm{G})\) such that every element of \(\mathbf{V}\) is the product of \(p\) commutators of elements of \(\mathbf{G}\) . (Let \(x_{1},y_{1},\ldots ,x_{p},y_{p}\) be elements of \(\mathbf{L}(\mathbf{G})\) such that the \([x_1,y_1]\) form a basis of \(\mathbf{L}(\mathbf{G})\) . For \(s,t\) in \(\mathbf{R}\) , we write

\[
\rho_ {i} (s, t) = (\exp s y _ {i}) ^ {- 1} (\exp t x _ {i}) ^ {- 1} (\exp s y _ {i}) (\exp t x _ {i}).
\]

Apply the implicit function theorem to the mapping

\[
\left(s _ {1}, t _ {1}, \dots , s _ {p}, t _ {p}\right) \mapsto \rho_ {1} \left(s _ {1}, t _ {1}\right) \dots \rho_ {p} \left(s _ {p}, t _ {p}\right)
\]

of \(\mathbf{R}^{2p}\) into (G, G).

\begin{enumerate}
\def\labelenumi{\arabic{enumi}.}
\setcounter{enumi}{10}
\item
  Let G be the semi-direct product of B = Z/2Z by K corresponding to the automorphism \(x \mapsto -x\) of the Lie group K. Then L(K) is an ideal of L(G) and L(B) = \{0\}, but (K, B) = K.
\item
  Let \((e_1, e_2, e_3)\) be the canonical basis of \(\mathbf{K}^3\) . Consider the nilpotent Lie algebra structure on \(\mathbf{K}^3\) such that \([e_1, e_2] = e_3\) , \([e_1, e_3] = [e_2, e_3] = 0\) . With the associated group law on \(\mathbf{K}^3\) (no. 5),
\end{enumerate}

\[
(x, y, z) (x ^ {\prime}, y ^ {\prime}, z ^ {\prime}) = (x + x ^ {\prime}, y + y ^ {\prime}, z + z ^ {\prime} + \frac {1}{2} (x y ^ {\prime} - y x ^ {\prime})).
\]

The mapping

\[
(\lambda , \mu , \nu) \mapsto (\exp \lambda e _ {1}) (\exp \mu e _ {2}) (\exp \nu (e _ {1} + e _ {3}))
\]

of \(K^{3}\) into G is neither surjective (show that \((0,1,1)\) is not in the image) nor injective (show that \((0,1,0)\) and \((1,1,-1)\) have the same image).

\begin{enumerate}
\def\labelenumi{\arabic{enumi}.}
\setcounter{enumi}{12}
\tightlist
\item
  \begin{enumerate}
  \def\labelenumii{(\alph{enumii})}
  \tightlist
  \item
    A multiplication is defined on \(\mathbf{R}^3\) as follows:
  \end{enumerate}
\end{enumerate}

\[
(x, y, z). (x ^ {\prime}, y ^ {\prime}, z ^ {\prime}) = (x + x ^ {\prime} \cos z - y ^ {\prime} \sin z, y + x ^ {\prime} \sin z + y ^ {\prime} \cos z, z + z ^ {\prime}).
\]

Show that a solvable real Lie group \(\mathbf{G}\) is thus obtained such that \(\mathrm{DG} = \mathbf{R}^2\times \{0\}\) . The centre \(\mathbf{Z}\) of \(\mathbf{G}\) is \(\{0\} \times \{0\} \times 2\pi \mathbf{Z}\) .

\begin{enumerate}
\def\labelenumi{(\alph{enumi})}
\setcounter{enumi}{1}
\tightlist
\item
  For \((x,y,z)\in G\) , let \(\pi (x,y,z)\) be the mapping
\end{enumerate}

\[
(\lambda , \mu) \mapsto (\lambda \cos z - \mu \sin z + x, \lambda \sin z + \mu \cos z + y)
\]

of \(\mathbf{R}^2\) into \(\mathbf{R}^2\) . Show that \(\pi\) is a morphism of \(G\) onto a Lie subgroup \(\mathbf{G}'\) of the affine group of \(\mathbf{R}^2\) , generated by the translations and rotations of \(\mathbf{R}^2\) . Show that \(\operatorname{Ker} \pi = \mathbf{Z}\) .

\begin{enumerate}
\def\labelenumi{(\alph{enumi})}
\setcounter{enumi}{2}
\item
  We canonically identify \(\mathrm{L}(\mathrm{G}) = \mathrm{T}_{(0,0,0)}\mathrm{G}\) with \(\mathbf{R}^3\) . Let \((e_1, e_2, e_3)\) be the canonical basis of \(\mathbf{R}^3\) . Show that \([e_1, e_2] = 0\) , \([e_3, e_1] = e_2\) , \([e_3, e_2] = -e_1\) .
\item
  Show that, for \(c \neq 0\) ,
\end{enumerate}

\[
\exp_ {\mathrm{G}} (a, b, c) = \left(\frac {1}{c} (a \sin c + b \cos c - b), \frac {1}{c} (- a \cos c + b \sin c + a), c\right).
\]

Deduce that, for all \(u \in L(G)\) , \(Z \subset \exp(\mathbf{R}u)\) . Show that \(\exp_G: L(G) \to G\) is neither injective nor surjective.

\begin{enumerate}
\def\labelenumi{\arabic{enumi}.}
\setcounter{enumi}{13}
\tightlist
\item
  \begin{enumerate}
  \def\labelenumii{(\alph{enumii})}
  \tightlist
  \item
    Let \(\mathbf{G} = \mathbf{GL}(n, \mathbf{C})\) . Show that \(\exp_{\mathbf{G}}\) is surjective. (Use the holomorphic functional calculus of Spectral Theories, Chapter I, § 4, no. 8.)
  \end{enumerate}
\end{enumerate}

\begin{enumerate}
\def\labelenumi{(\alph{enumi})}
\setcounter{enumi}{1}
\tightlist
\item
  Let \(\mathbf{G}' = \mathbf{SL}(2, \mathbf{C})\) . Show that, if \(\sigma \in \mathbf{C}^*\) , the element \(\left( \begin{array}{cc} -1 & \sigma \\ 0 & -1 \end{array} \right)\)
\end{enumerate}

\(G'\) does not belong to the image of \(\exp_{G'}\) .

\begin{enumerate}
\def\labelenumi{\arabic{enumi}.}
\setcounter{enumi}{14}
\tightlist
\item
  \begin{enumerate}
  \def\labelenumii{(\alph{enumii})}
  \tightlist
  \item
    Let \(x \in \mathfrak{sl}(2, \mathbf{R})\) and \(\Delta = \det x\) . Show that
  \end{enumerate}
\end{enumerate}

\[
e ^ {x} = \operatorname{ch} \sqrt {- \Delta}. I + \frac {\operatorname{sh} \sqrt {- \Delta}}{\sqrt {- \Delta}} \cdot x \quad \mathrm{if} \Delta <   0
\]

\[
e ^ {x} = \cos {\sqrt {\Delta}}. I + \frac {\sin {\sqrt {\Delta}}}{\sqrt {\Delta}}. x \qquad \mathrm{if} \Delta > 0.
\]

\begin{enumerate}
\def\labelenumi{(\alph{enumi})}
\setcounter{enumi}{1}
\tightlist
\item
  Let \(\mathbf{H} = \mathbf{SL}(2, \mathbf{R})\) . Show that \(g = \begin{pmatrix} \lambda & 0 \\ 0 & \lambda^{-1} \end{pmatrix} \in \mathbf{H}\) is the image of \(\exp_{\mathbf{H}}\) if and only if \(\lambda > 0\) or \(\lambda = -1\) . If \(\lambda = 0\) , all the one-parameter subgroups containing \(g\) are equal.
\end{enumerate}

\begin{enumerate}
\def\labelenumi{\arabic{enumi}.}
\setcounter{enumi}{15}
\tightlist
\item
  Let G be a finite-dimensional Lie group. Show that there exists a basis \((x_{1},\ldots,x_{n})\) of L(G) such that, for all i, \(\exp(\mathbf{R}x_{i})\) is a Lie subgroup of G. (Use Spectral Theories, Chapter II, §2, Lemma 1.)
\end{enumerate}

¶ 17. (a) Let c denote a real solvable Lie algebra with a basis \((a, b, c)\) such that \([a, b] = c\) , \([a, c] = -b\) , \([b, c] = 0\) . Let d denote a real solvable Lie algebra with a basis \((a, b, c, d)\) such that \([a, b] = c\) , \([a, c] = -b\) , \([b, c] = d\) , \([a, d] = [b, d] = [c, d] = 0\) . Let g be a real solvable Lie algebra. If g contains non-zero elements x, y, z such that \([x, y] = z\) and \([x, z] = -y\) , then g contains a subalgebra isomorphic to c or d.~(Write \(a_{1} = y\) , \(b_{1} = z\) , \(c_{1} = [y, z]\) ; define \(a_{i}\) , \(b_{i}\) , \(c_{i}\) inductively by \(a_{i} = [a_{i-1}, c_{i-1}]\) , \(b_{i} = [b_{i-1}, c_{i-1}]\) , \(c_{i} = [a_{i}, b_{i}]\) . Consider the smallest integer k such that \(c_{k} = 0\) .)

\begin{enumerate}
\def\labelenumi{(\alph{enumi})}
\setcounter{enumi}{1}
\item
  Let g and \(g^{*}\) be real solvable Lie algebras and \(\phi\) a homomorphism of g onto \(g^{*}\) . If \(g^{*}\) contains a subalgebra isomorphic to c or d, g has the same property.
\item
  Let \(\mathfrak{h}\) be a complex solvable Lie algebra. Consider a Jordan-Hölder series for the adjoint representation of \(\mathfrak{h}\) ; the quotients of this series define 1-dimensional representations of \(\mathfrak{h}\) and hence linear forms on \(\mathfrak{h}\) . These linear forms, which depend only on \(\mathfrak{h}\) , are called the roots of \(\mathfrak{h}\) . If \(\mathfrak{h}'\) is a real solvable Lie algebra, the roots of \(\mathfrak{h}'\) are the restrictions to \(\mathfrak{h}'\) of the roots of \(\mathfrak{h}' \otimes_{\mathbf{R}} \mathbf{C}\) .
\end{enumerate}

Then let G be a finite-dimensional simply connected solvable real Lie group and g its Lie algebra. Show that the following conditions are equivalent: (α) exp₄ is injective; (β) exp₅ is surjective; (γ) exp₆ is bijective; (δ) exp₇ is an isomorphism of the analytic manifold L(G) onto the analytic manifold G; (ε) L(G) contains no subalgebra isomorphic to c or d; (ζ) there exists no quotient algebra of L(G) admitting a subalgebra isomorphic to c; (η) every root of g is of the form \(\phi + i\delta'\) where \(\phi\) , \(\phi'\) are in \(g^*\) and \(\phi'\) is proportional to \(\phi\) ; (θ) for all \(x \in g\) , the only pure imaginary eigenvalue (in C) of ad x is 0.†
% semantic-fragments: legacy-input-seam
\relax{}

\begin{enumerate}
\def\labelenumi{\arabic{enumi}.}
\setcounter{enumi}{17}
\tightlist
\item
  Let \(G\) be a connected real Lie group, \(Z\) its centre, \(Z_0\) the identity component of \(Z\) , \(g\) the Lie algebra of \(G\) and \(h\) the centre of \(g\) . Then \(Z\) and \(Z_0\) are Lie quasi-subgroups of \(G\) with Lie algebra \(h\) (Exercise 3). An acceptable norm on \(g\) is a norm defining the topology of \(g\) and making \(g\) into a normed Lie algebra.
\end{enumerate}

\begin{enumerate}
\def\labelenumi{(\alph{enumi})}
\item
  For all \(r > 0\) and \(z \in \mathfrak{z}\) , there exists an acceptable norm on \(\mathfrak{g}\) whose value at \(z\) is \(< r\) . (Let \(q\) be an acceptable norm on \(\mathfrak{g}\) . Show that, for suitably chosen \(\lambda > 0\) , the function \(x \mapsto \|x\| = \lambda q(x) + \inf_{y \in \mathfrak{z}} q(x + y)\) has the required properties.)
\item
  Show that the following conditions are equivalent:
\item
  \(Z_{0}\) is simply connected;
\end{enumerate}

\begin{enumerate}
\def\labelenumi{(\roman{enumi})}
\setcounter{enumi}{1}
\item
  for every acceptable norm on g, the restriction of \(\exp_{G}\) to the open ball of centre 0 and radius \(\pi\) is injective;
\item
  there exists \(r > 0\) such that, for every acceptable norm on \(\mathfrak{g}\) , the restriction of \(\exp_{\mathbb{G}}\) to the open ball of centre 0 and radius \(r\) is injective.
\end{enumerate}

(To prove (iii) \(\Rightarrow\) (i), use (a). To prove (i) \(\Rightarrow\) (ii), suppose that \(x, y\) are in \(\mathfrak{g}\) , \(\| x \| < \pi\) , \(\| y \| < \pi\) , \(x \neq y\) , \(\exp_{\mathfrak{g}} x = \exp_{\mathfrak{g}} y\) . Then \(\exp \mathrm{d} x = \exp \mathrm{d} y\) and hence \(\mathrm{ad} x = \mathrm{ad} y\) by Proposition 17 of §6 applied to a complexification of \(\mathfrak{g}\) . Hence there exists a non-zero \(z\) in \(\mathfrak{z}\) such that \(\exp_{\mathfrak{g}} z = e\) ; it follows that (i) is false.)

\begin{enumerate}
\def\labelenumi{(\alph{enumi})}
\setcounter{enumi}{2}
\tightlist
\item
  By considering the group \(\mathbf{G}'\) of Exercise 13 (b), show that the conclusion of (b) is no longer true if \(\pi\) is replaced by a number \(> \pi\) . (In the notation of Exercise 13, use the norm \(ae_{1} + be_{2} + ce_{3} \mapsto (a^{2} + b^{2} + c^{2})^{\frac{1}{2}}\) on \(\mathrm{L}(\mathbf{G}')\) .)
\end{enumerate}

\begin{enumerate}
\def\labelenumi{\arabic{enumi}.}
\setcounter{enumi}{18}
\item
  Let G be a locally compact group and H a closed subgroup of G. Suppose that H is a finite-dimensional simply connected solvable Lie group and that G/H is compact. Show that there exists a compact subgroup L of G such that G is the semi-direct product of L and H. (Argue by induction on dim H. Consider the last non-trivial derived group of H and use Integration, Chapter VII, § 3, Proposition 3.)
\item
  Let \(G\) be a finite-dimensional solvable connected real Lie group. We introduce the notation \(S, S', F, \sigma\) of §6, no. 10, proof of Proposition 20.
\end{enumerate}

\begin{enumerate}
\def\labelenumi{(\alph{enumi})}
\item
  Using Proposition 21, show that \(\sigma\) is an isomorphism of S onto a Lie subgroup of the underlying real Lie group of S'.
\item
  Deduce that the universal complexification \(\tilde{G}\) of \(G\) is identified with \(S' / \sigma(F)\) and that the canonical mapping of \(G\) into \(\tilde{G}\) is an isomorphism of \(G\) onto a Lie subgroup of the underlying real Lie group of \(\tilde{G}\) .
\end{enumerate}

¶ 21. (a) Let G be a simply connected solvable real Lie group with the following properties: (α) L(G), which is n-dimensional, has a commutative ideal of dimension n - 1, corresponding to a subgroup A of G; (β) there exists an element σ of the centre of G which does not belong to A. Show that there exists an element \(x\) of \(\mathrm{L}(\mathbf{G})\) such that \(\exp x = \sigma\) . Show that \(\mathrm{L}(\mathbf{G})\) is the product of a commutative ideal and an ideal with a basis

\[
(x, a _ {1}, b _ {1}, \dots , a _ {k}, b _ {k})
\]

such that \(a_1, b_1, \ldots, a_k, b_k\) belong to \(\mathrm{L}(\mathrm{A})\) , \([x, a_i] = 2\pi n_i b_i\) , \([x, b_i] = -2\pi n_i a_i\) ( \(n_i \in \mathbf{Z} - \{0\}\) ) for all \(i\) . Generalize the results of Exercise 13 (d) to G.

\begin{enumerate}
\def\labelenumi{(\alph{enumi})}
\setcounter{enumi}{1}
\tightlist
\item
  Let G be a finite-dimensional simply connected solvable real Lie group. Let D be a discrete subgroup of the centre of G. There exist a basis \((x_{1}, x_{2}, \ldots, x_{n})\) of L(G) and an integer \(r \leqslant n\) with the following properties: (α) every element of G can be written uniquely in the form
\end{enumerate}

\[
\left(\exp t _ {1} x _ {1}\right) \dots \left(\exp t _ {n} x _ {n}\right)
\]

where \(t_1, \ldots, t_n\) are in \(\mathbf{R}\) ; \((\beta)x_1, \ldots, x_r\) are pairwise permutable and

\[
(\exp x _ {1}, \dots , \exp x _ {r})
\]

is a basis of the commutative group D. (Argue by induction on the dimension of G. Let a be maximal commutative ideal of L(G). Let A be the corresponding integral subgroup. Then A is closed in G and DA/A is a discrete subgroup of the centre of G/A, to which the induction hypothesis can be applied, whence there are elements \(x_1^*, \ldots, x_m^*\) of L(G/A) and an integer \(s\) . For \(1 \leqslant i \leqslant s\) , let \(\sigma_i\) be an element of D whose class modulo A is exp \(x_1^*\) . Construct one by one, using (a), representatives \(x_1, \ldots, x_m\) of \(x_1^*, \ldots, x_m^*\) such that exp \(x_i = \sigma_i\) and \([x_1, x_j] = 0\) for \(1 \leqslant i, j \leqslant s\) .

\begin{enumerate}
\def\labelenumi{(\alph{enumi})}
\setcounter{enumi}{2}
\tightlist
\item
  Deduce from (b) that every finite-dimensional connected solvable real Lie group is homomorphic to a space \(\mathbf{R}^n\times \mathbf{T}^m\) ( \(m,n\) integers \(\geqslant 0\) ).
\end{enumerate}

\begin{enumerate}
\def\labelenumi{\arabic{enumi}.}
\setcounter{enumi}{21}
\item
  Let \(G\) be a finite-dimensional connected real Lie group such that \(L(G)\) is reductive. Then \(G = (V \times S)/N\) where \(V\) is a finite-dimensional real vector space, where \(S\) is a simply connected semi-simple real Lie group and \(N\) is a central discrete subgroup of \(V \times S\) . Then \(G/\overline{D^1}G\) is isomorphic to \(V/\mathrm{pr}_1\overline{N}\) . Hence \(G/\overline{D^1}N\) is compact if and only if \(\mathrm{pr}_1\mathrm{N}\) generates the vector space \(V\) .
\item
  \begin{enumerate}
  \def\labelenumii{(\alph{enumii})}
  \tightlist
  \item
    Let G be a finite-dimensional commutative real Lie group with only a finite number of connected components. For G to be compact, it is necessary and sufficient that every finite-dimensional analytic linear representation of G on a complex vector space be semi-simple. (Use the proof of Proposition 32.)
  \end{enumerate}
\end{enumerate}

\begin{enumerate}
\def\labelenumi{(\alph{enumi})}
\setcounter{enumi}{1}
\tightlist
\item
  Let G be a finite-dimensional commutative complex Lie group with only a finite number of connected components. Let N be the kernel of \(\exp_{G}\) .
\end{enumerate}

For N to generate over C the vector space \(L(G)\) , it is necessary and sufficient that every finite-dimensional analytic linear representation of G be semi-simple. (Use (a), Lemma 1 and Proposition 32.)

\begin{enumerate}
\def\labelenumi{\arabic{enumi}.}
\setcounter{enumi}{23}
\item
  The Lie group \(\mathbf{SL}(2,\mathbf{R})\) is connected and almost simple, but \(\{I,-I\}\) is a normal commutative subgroup of \(\mathbf{SL}(2,\mathbf{R})\) .
\item
  Let G be an almost simple connected real or complex Lie group. Let A be a normal subgroup of G. If \(A \neq G\) , A is discrete and central. (Use Exercise 8 of §4.) Therefore the quotient of G by its centre is simple as an abstract group.
\item
  Let G be a finite-dimensional connected real Lie group. Suppose that G admits a finite-dimensional injective continuous linear representation \(\rho\) . Then (G, G) is closed in G. (Let R be the radical of G and S a maximal semi-simple integral subgroup of G. Using Exercise 7 (b), reduce it to the case where G is the semi-direct product of S and R. By Chapter I, § 6, Proposition 6, \(\rho\) is unipotent on (G, R). Hence \(\rho((G, R))\) is closed in the linear group and therefore (G, R) is closed in G.)
\item
  Let G be a finite-dimensional solvable simply connected real Lie group and N the largest connected nilpotent normal subgroup of G. Then G admits a finite-dimensional injective continuous linear representation which is unipotent on N. (Argue by induction on the dimension of G. Use Proposition 20 and Chapter I, § 7, Theorem 1.)†
\end{enumerate}

¶ 28. (a) Let k be a commutative field of characteristic 0, L an n-dimensional Lie algebra over k, \(L_{1}\) and \((n-1)\) -dimensional subalgebra containing no non-zero ideal of L and \(a_{0}\) an element of L not belonging to \(L_{1}\) . For \(i=2,3,\ldots\) , we define \(L_{i}\) inductively to be the set of \(x\in L_{i-1}\) such that \([x,a_{0}]\in L_{i-1}\) . We write \(L_{i}=L\) for \(i\leqslant0\) . Show that \([L_{0},L_{i}]\subset L_{i+j}\) (by induction on \(i+j\) ) and then that, for \(0\leqslant i\leqslant n\) , \(L_{i}\) is a subalgebra of L of codimension i (by induction on i).

\begin{enumerate}
\def\labelenumi{(\alph{enumi})}
\setcounter{enumi}{1}
\tightlist
\item
  For 0 \textless{} i \textless{} n, choose in \(L_{i}\) an element \(a_{i}\) not belonging to \(L_{i+1}\) such that \([a_{0}, a_{i}] \equiv ia_{i-1} \pmod{L_{i}}\) . Show, by induction on the ordered pairs \((i, j)\) ordered lexicographically, that, for \(0 \leqslant i < j < n\) , \(i + j - 1 < n\) and
\end{enumerate}

\[
[ a _ {i}, a _ {j} ] \equiv (j - i) a _ {i + j - 1} (\mathrm{mod}. \mathrm{L} _ {i + j}).
\]

\begin{enumerate}
\def\labelenumi{(\alph{enumi})}
\setcounter{enumi}{2}
\item
  Deduce that L is either of dimension 1, or of dimension 2 and noncommutative, or isomorphic to \(\mathfrak{sl}(2,k)\) .
\item
  Let A be a finite-dimensional real Lie algebra. A subalgebra B of A is called extendable if there exists a subalgebra \(B_{1} \supset B\) such that dim \(B_{1} = \dim B + 1\) . Let \(R'(A)\) denote the intersection of all the nonextendable subalgebras of A. Let \(R(A)\) denote the largest ideal of A with a composition series as an A-module all of whose quotients are of dimension 1.
\end{enumerate}

Show that \(\mathbf{R}'(\mathbf{A})\) is a characteristic ideal of \(\mathbf{A}\) . (Observe that \(\mathbf{R}'(\mathbf{A})\) is stable under \(\operatorname{Aut}(\mathbf{A})\) .)

\begin{enumerate}
\def\labelenumi{(\alph{enumi})}
\setcounter{enumi}{4}
\item
  Show that \(\mathbf{R}'(\mathbf{A})\supset \mathbf{R}(\mathbf{A})\) and that \(\mathrm{R}'(\mathrm{A} / \mathrm{R}(\mathrm{A})) = \mathrm{R}'(\mathrm{A}) / \mathrm{R}(\mathrm{A})\)
\item
  Show that, if B is a subalgebra of A, then R'(B) ⊃ R'(A) ∩ B.
\item
  If A is solvable, \(\mathrm{R}^{\prime}(\mathrm{A}) = \mathrm{R}(\mathrm{A})\) . (It can be assumed that \(\mathrm{R}(\mathrm{A}) = \{0\}\) , by (e). Suppose that \(\mathrm{R}^{\prime}(\mathrm{A}) \neq \{0\}\) . By (d), there exists a minimal non-zero ideal I of A contained in \(\mathrm{R}^{\prime}(\mathrm{A})\) . Using (f) and Chapter I, §5, Corollary 1 to Theorem 1, show that dim I = 1, whence a contradiction.)
\item
  Let P be the radical of \(\mathrm{R}^{\prime}(\mathrm{A})\) , B a semi-simple subalgebra of A and \(C = P + B\) which is a subalgebra of A by (d). Using (f), show that \(ad_{P}B\) can be expressed in a triangular form with respect to a suitable basis of P. Conclude that \([P, B] = \{0\}\) .
\end{enumerate}

\((k)\) R(A) is the radical of \(\mathrm{R}'(\mathrm{A})\) . (Use \((e)\) , \((f)\) , \((g)\) , \((h)\) and a Levi decomposition of A.) Let \(\mathrm{R}'(\mathrm{A}) = \mathrm{R}(\mathrm{A}) \oplus \mathrm{T}\) be a Levi decomposition of \(\mathrm{R}'(\mathrm{A})\) . By \((h)\) , \(\mathrm{R}'(\mathrm{A}) = \mathrm{R}(\mathrm{A}) \times \mathrm{T}\) . Then \(\mathrm{R}(\mathrm{T}) = \{0\}\) . Applying \((c)\) to the simple factors of T, conclude that \(\mathrm{T} = \{0\}\) . It has therefore been shown that \(\mathrm{R}(\mathrm{A}) = \mathrm{R}'(\mathrm{A})\) .

\begin{enumerate}
\def\labelenumi{(\alph{enumi})}
\setcounter{enumi}{11}
\item
  Let G be a connected real Lie group with Lie algebra A. Let \(\mathcal{S}(G)\) be the set of subgroups N of G such that there exists a decreasing sequence \((\mathrm{N}_m, \mathrm{N}_{m-1}, \ldots, \mathrm{N}_0)\) of subgroups with the following properties: \(\mathrm{N}_m = \mathrm{N}\) , \(\mathrm{N}_0 = \{\varepsilon\}\) , each \(\mathrm{N}_i\) is a normal connected Lie subgroup of G and \(\dim \mathrm{N}_i / \mathrm{N}_{i-1} = 1\) for all \(i > 0\) . Show that the integral subgroup R(G) of G with Lie algebra R(A) is the largest element of \(\mathcal{S}(G)\) . (Argue by induction on \(\dim \mathrm{R(A)}\) . Let I be a 1-dimensional ideal of A and N the corresponding integral subgroup of G. Pass to the quotient by N. If N is not closed, use Exercise 21 (a) of § 6.)
\item
  A Lie subgroup H of G is called extendable if there exists a Lie subgroup \(H_{1} \supset H\) such that \(\dim H_{1} = \dim H + 1\) . Let \(R'(G)\) denote the intersection of all the non-extendable connected Lie subgroups of G.
\end{enumerate}

Let B be a non-extendable subalgebra of A. Show that the corresponding integral subgroup of G is closed. (Use Proposition 5.) Deduce that

\[
\mathrm{L} (\mathrm{R} ^ {\prime} (\mathrm{G})) \subset \mathrm{R} (\mathrm{A}).
\]

whence \(\mathrm{R}^{\prime}(\mathrm{G})\subset\mathrm{R}(\mathrm{G})\) .

\begin{enumerate}
\def\labelenumi{(\alph{enumi})}
\setcounter{enumi}{13}
\tightlist
\item
  Show that \(\mathrm{R}(\mathrm{G}) = \mathrm{R}'(\mathrm{G})\) . (Reduce it to the case where \(\mathrm{R}'(\mathrm{G}) = \{e\}\) . Then use Exercise 22 of §6.)†
\end{enumerate}

\begin{enumerate}
\def\labelenumi{\arabic{enumi}.}
\setcounter{enumi}{28}
\tightlist
\item
  A real Lie group G is said to be of type (N) if it is finite-dimensional,
\end{enumerate}

nilpotent, connected and simply connected. If G is such a group and g is its Lie algebra, the mapping exp: g → G is an isomorphism when g is given the group structure defined by the Hausdorff law (cf.~Chapter II, §6, no. 5, Remark 3). Let log: G → g denote the inverse isomorphism.

\begin{enumerate}
\def\labelenumi{(\alph{enumi})}
\item
  Let V be a vector Q-subspace of g. Show the equivalence of the following conditions:
\item
  V is a Lie \(\mathbf{Q}\) -subalgebra of \(\mathfrak{g}\) ;
\end{enumerate}

\begin{enumerate}
\def\labelenumi{(\roman{enumi})}
\setcounter{enumi}{1}
\tightlist
\item
  \(\exp(V)\) is a subgroup of G.
\end{enumerate}

(Use Exercise 5 of § 6 of Chapter II.)

For a subgroup H of G to be of the form \(\exp(V)\) , where V is a vector Q-subspace of g, it is necessary and sufficient that H be saturated in G (Chapter II, §4, Exercise 14), i.e.~that the relations \(x \in G\) , \(x^{n} \in H\) , \(n \neq 0\) imply \(x \in H\) (loc. cit.). If this is so, show that H is an integral subgroup of G if and only if \(\log(H)\) is a Lie R-subalgebra of g.

\begin{enumerate}
\def\labelenumi{(\alph{enumi})}
\setcounter{enumi}{1}
\item
  Let V be a Lie Q-subalgebra of g of finite dimension m, let \((e_{1},\ldots,e_{m})\) be a basis of V over Q and let \(\Lambda\) be the subgroup of V generated by \(e_{1},\ldots,e_{m}\) . Using the fact that the Hausdorff law is polynomial, show that there exists an integer \(d\geqslant1\) such that, for every integer r which is a non-zero multiple of d, \(\exp(r\Lambda)\) is a subgroup of G. Let d be such an integer and r a multiple of d.~Show that \(\exp(r\Lambda)\) is discrete if and only if \((e_{1},\ldots,e_{m})\) is a free family over R, i.e.~if the canonical mapping of V \(\otimes_{q}R\) into g is injective. Suppose that this is the case; show that G/exp(r \(\Lambda\) ) is compact if and only if V \(\otimes_{q}R\to g\) is bijective, i.e.~if V is a Q-form of g. (To show that the condition is sufficient, argue by induction on the nilpotency class of g.)
\item
  Conversely, let \(\Gamma\) be a discrete subgroup of G, let \(\Gamma\) be its saturation in G (Chapter II, loc. cit.) and let \(\mathfrak{g}_{\Gamma} = \log (\Gamma) = \mathbf{Q}. \log (\Gamma)\) be the corresponding Lie \(\mathbf{Q}\) -subalgebra. Suppose that \(\mathbf{R}, \mathfrak{g} = \mathfrak{g}\) , i.e.~that \(\Gamma\) is contained in no distinct integral subgroup of G. Let \(z\) be an element \(\neq 1\) of the centre of \(\Gamma\) and let \(x = \log (z)\) . Show that \(x\) belongs to the centre of \(\mathfrak{g}\) . If \(\mathrm{X} = \exp (\mathbf{R}x)\) , show that \(\mathrm{X} / (\Gamma \cap \mathrm{X})\) is compact and that the image of \(\Gamma\) in \(\mathrm{G} / \mathrm{X}\) is discrete. Deduce, arguing by induction on dim G, that \(\mathrm{G} / \Gamma\) is compact and that \(\mathfrak{g}_{\Gamma}\) is a G-form of \(\mathfrak{g}\) . If \(\Lambda\) is a lattice of \(\mathfrak{g}_{\Gamma}\) , show that there exists an integer \(d \neq 0\) such that \(\Gamma\) contains \(\exp (d\Lambda)\) and that the index of \(\exp (d\Lambda)\) in \(\Gamma\) is then finite.
\end{enumerate}

Let H be an integral subgroup of G with Lie algebra \(\mathfrak{h}\) . Show that \(\mathrm{H} / (\mathrm{H} \cap \Gamma)\) is compact if and only if \(\mathfrak{h}\) is rational with respect to the \(\mathbf{Q}\) -structure \(g_{\Gamma}\) (in particular if \(\mathfrak{h}\) is one of the terms of the lower---or upper---central series of \(g\) ).

\begin{enumerate}
\def\labelenumi{(\alph{enumi})}
\setcounter{enumi}{3}
\item
  Let \(\Gamma\) be a discrete subgroup of \(G\) , let \(H\) be the smallest integral subgroup of \(G\) containing \(\Gamma\) and let \(\mathfrak{h}\) be its Lie algebra. Show that \(H / \Gamma\) is compact and that \(\mathfrak{g} \otimes_{\mathbf{q}} \mathbf{R} \to \mathfrak{g}\) is injective and has image \(\mathfrak{h}\) . (Apply (c) to the nilpotent group \(H\) .)
\item
  Let \(\Gamma\) be a discrete subgroup of \(G\) . Show that there exists a basis \((x_{1},\ldots ,x_{q})\) of \(L(G)\) with the following properties:
\item
  for all \(i \in [1, q]\) , \(\mathbf{R}x_i + \cdots + \mathbf{R}x_q\) is an ideal \(\mathfrak{n}_i\) of \(\mathbf{R}x_{i-1} + \cdots + \mathbf{R}x_q\) ;
\end{enumerate}

\begin{enumerate}
\def\labelenumi{(\roman{enumi})}
\setcounter{enumi}{1}
\tightlist
\item
  there exists \(p \in (1, q]\) such that \(\Gamma\) is the set of products
\end{enumerate}

\[
\exp (m _ {p} x _ {p}) \exp (m _ {p + 1} x _ {p + 1}) \dots \exp (m _ {q} x _ {q})
\]

where \(m_{p}, \ldots, m_{q}\) are in Z;

\begin{enumerate}
\def\labelenumi{(\roman{enumi})}
\setcounter{enumi}{2}
\tightlist
\item
  if \(G / \Gamma\) is compact, \(\{n_1, n_2, \ldots, n_q\}\) contains the lower central series of \(g\) . (Using (d) and Proposition 16, reduce it to the case where \(G / \Gamma\) is compact. Then argue by induction on \(\dim G\) , using (c).)
\end{enumerate}

\begin{enumerate}
\def\labelenumi{(\arabic{enumi})}
\setcounter{enumi}{29}
\tightlist
\item
  Give an example of a 7-dimensional real nilpotent Lie algebra with no basis with respect to which the constants of structure are rational (cf.~Chapter I, § 4, Exercise 18). Deduce that the corresponding group of type (N) (Exercise 29) has no discrete subgroup with compact subgroup. (Use Exercise 29.)
\end{enumerate}

\begin{enumerate}
\def\labelenumi{\arabic{enumi}.}
\setcounter{enumi}{30}
\item
  Let G and \(G'\) be two Lie groups of type (N) (Exercise 29) and let \(\Gamma\) be a discrete subgroup of G such that \(G/\Gamma\) is compact. Show that every homomorphism \(f: \Gamma \to G'\) can be extended uniquely to a Lie group morphism of G into \(G'\) (begin by extending f to the saturation \(\tilde{\Gamma}\) of \(\Gamma\) and derive a homomorphism of the Lie Q-algebra \(\log(\tilde{\Gamma})\) into the Lie algebra of \(G'\) , cf.~Exercise 29).
\item
  Let G be a Lie group of type (N) (Exercise 29) and let \(\Gamma\) be a discrete subgroup of G. Prove the equivalence of the following conditions:
\end{enumerate}

\begin{enumerate}
\def\labelenumi{(\alph{enumi})}
\item
  G/Γ is compact;
\item
  the measure of \(\mathrm{G} / \Gamma\) is finite (relative to a non-zero G-invariant positive measure);
\item
  every integral subgroup of G containing \(\Gamma\) is equal to G. (Use Exercise 29.)
\end{enumerate}

¶ 33. Let \(\Gamma\) be a group. Prove the equivalence of the following conditions: (a) \(\Gamma\) is nilpotent, torsion-free and finitely generated;

\begin{enumerate}
\def\labelenumi{(\alph{enumi})}
\setcounter{enumi}{1}
\item
  there exists a Lie group of type (N) (Exercise 29) containing \(\Gamma\) as a discrete subgroup;
\item
  there exists a Lie group G of type (N) (Exercise 29) containing \(\Gamma\) as a discrete subgroup and such that \(G / \Gamma\) is compact.
\end{enumerate}

(The equivalence of \(\langle b\rangle\) and \(\langle c\rangle\) follows from Exercise 29. The implication \((\varepsilon)\Rightarrow \langle a\rangle\) is proved by induction on \(\dim (\mathrm{G})\) by the method of Exercise 29 \((c)\) . To prove that \((a)\Rightarrow (c)\) , show first that the Lie \(\mathbf{Q}\) -algebra attached to the saturation of \(\Gamma\) (cf.~Chapter II, § 6, Exercise 4) is finite-dimensional; then take the tensor product of this algebra with \(\mathbf{R}\) and the corresponding group of type (N.)

\begin{enumerate}
\def\labelenumi{\arabic{enumi}.}
\setcounter{enumi}{33}
\tightlist
\item
  Let \(G\) be a finite-dimensional connected real or complex Lie group and \(\rho\) an analytic linear representation of \(G\) on a finite-dimensional complex vector space \(V\) . Suppose that \(\rho\) is semi-simple. The group \(G\) operates by automorphisms on the algebra \(S(V^*)\) of polynomial functions on \(V\) .
\end{enumerate}

\begin{enumerate}
\def\labelenumi{(\alph{enumi})}
\item
  Let \(\mathsf{S}(\mathbf{V}^{*})^{\mathbb{G}}\) be the set of G-invariant elements of \(\mathsf{S}(\mathbf{V}^{*})\) . There exists a projector \(p\) of \(S(V^{*})\) onto \(S(V^{*})^{\mathbb{G}}\) which commutes with the operations of \(G\) and which leaves stable every \(G\) -stable vector subspace of \(S(V^{*})\) .
\item
  Let I, J be ideals of \(S(V^*)\) which are stable under G. Let A (resp. B) be the set of zeros of I (resp. J) in V. Suppose that \(A \cap B = \varnothing\) . There then exists \(u \in I\) such that \(u\) is G-invariant and that \(u = 1\) in B. (By Commutative Algebra, Chapter 7, §3, no. 3, Proposition 2, there exists \(v \in I\) , \(w \in J\) such that \(v + w = 1\) Let \(u = pv\) . Show that \(u\) has the required properties.)
\end{enumerate}

\begin{enumerate}
\def\labelenumi{\arabic{enumi}.}
\setcounter{enumi}{34}
\tightlist
\item
  Let \(\mathbf{G}\) be a compact subgroup of \(\mathbf{GL}(n,\mathbf{R})\) .
\end{enumerate}

\begin{enumerate}
\def\labelenumi{(\alph{enumi})}
\item
  Let A, B be disjoint compact G-invariant subsets of \(\mathbf{R}^{\mathbf{n}}\) . There exists a G-invariant polynomial function \(u\) on \(\mathbf{R}^{\mathbf{n}}\) , such that \(u = 1\) on \(\mathrm{A}\) and \(u = 0\) on B. (There exists a continuous real-valued function \(v\) on \(\mathbf{R}^{\mathbf{n}}\) such that \(v = 1\) on A and \(v = 0\) on B. By the Stone-Weierstrass Theorem, there exists a polynomial function \(w\) on \(\mathbf{R}^{\mathbf{n}}\) such that \(|v - w| \leqslant \frac{1}{3}\) on \(\mathrm{A} \cup \mathrm{B}\) . Derive \(u\) by an integration process with respect to the normalized Haar measure of G.)
\item
  Let \(x \in \mathbf{R}^n\) . Then \(Gx\) is the set of zeros in \(\mathbf{R}^n\) of a finite number of G-invariant polynomials. (Use (a) and Corollary 2 to Theorem 1, no. 8.)
\end{enumerate}

\begin{enumerate}
\def\labelenumi{\arabic{enumi}.}
\setcounter{enumi}{35}
\tightlist
\item
  In \(\mathbf{SL}(2, \mathbf{R})\) , the elements which can be expressed in the form \((a, b)\) , where \(a, b\) in \(\mathbf{SL}(2, \mathbf{R})\) , are the elements \(\neq -1\) .
\end{enumerate}

¶ 37. Let G be a compact real Lie group and G' a finite-dimensional connected real Lie group. Suppose that L(G) is simple. Let \(\rho: G \to G'\) be a homomorphism of abstract groups. Suppose that there exists a neighbourhood V of \(e_{\alpha}\) in G such that \(\rho(V)\) is relatively compact. Then \(\rho\) is continuous. (Let V' be a neighbourhood of \(e_{\alpha'}\) in G'. Let V'\,' \(\subset\) V' be a neighbourhood of \(e_{\alpha'}\) , such that

\[
(x \in \mathrm{V} ^ {\prime \prime}, x _ {j} \in \rho (\mathrm{V}), y _ {j} \in \rho (\mathrm{V}) \text {   for   } 1 \leqslant j \leqslant n) \Rightarrow \left(\prod_ {i = 1} ^ {n} x _ {i} (y _ {i}, x) x _ {i} ^ {- 1} \in \mathrm{V} ^ {\prime}\right)
\]

where \(n = \dim G\) . It can be assumed that \(\rho\) is non-trivial. Then \(\operatorname{Ker} \rho\) is finite by Exercise 25. Therefore \(\rho(G)\) is not countable and hence not discrete. Hence there exists \(g \in G\) with non-open centralizer such that \(\rho(g) \in V''\) . Using Exercise 9 of §4 and its notation, \(\rho(M(g, n, V)) \subset V'\) and \(M(g, n, V)\) is a neighbourhood of \(e_G\) in G.)

\begin{enumerate}
\def\labelenumi{\arabic{enumi}.}
\setcounter{enumi}{37}
\tightlist
\item
  Let \(G\) be a finite-dimensional real Lie group and \(G_0\) its identity component. Consider the following property:
\end{enumerate}

\begin{enumerate}
\def\labelenumi{(\Alph{enumi})}
\setcounter{enumi}{5}
\tightlist
\item
  Ad(G) is closed in Aut L(G).
\end{enumerate}

Show that G has property (F) in each of the following cases:

\begin{enumerate}
\def\labelenumi{(\roman{enumi})}
\item
  G is connected and nilpotent;
\item
  every derivation of L(G) is inner;
\item
  \(\mathrm{G}_0\) has property (F) and \(\mathrm{G} / \mathrm{G}_0\) is finite;
\item
  G is an upper triangular group;
\item
  \(G_{0}\) is semi-simple.
\end{enumerate}

\begin{enumerate}
\def\labelenumi{\arabic{enumi}.}
\setcounter{enumi}{38}
\tightlist
\item
  Let G be a finite-dimensional connected real Lie Group, H an integral subgroup of G and \(\bar{H}\) its closure in G. Suppose that H has property (F) (Exercise 38).
\end{enumerate}

\begin{enumerate}
\def\labelenumi{(\alph{enumi})}
\item
  Let \(x \in \overline{\mathbf{H}}\) . Then \(\mathrm{Ad}_{\mathrm{L(G)}}x\) leaves \(\mathrm{L(H)}\) stable (no. 2, Proposition. 5); let \(u(x)\) be its restriction to \(\mathrm{L(H)}\) . Then \(u(\overline{\mathrm{H}}) = \mathrm{Ad}_{\mathrm{L(H)}}(\mathrm{H})\) . (Observe that \(\mathrm{Ad}_{\mathrm{L(H)}}(\mathrm{H})\) is dense in \(u(\overline{\mathrm{H}})\) and apply property (F).)
\item
  Let C be the centre of \(\overline{\mathrm{H}}\) . Then \(\overline{\mathrm{H}} = \mathrm{C.H}\) and C is the closure of the centre of H. (By (a), if \(x \in \overline{\mathrm{H}}\) , there exists \(y \in \mathrm{H}\) such that \(\mathrm{Ad}_{\mathrm{L(H)}} x = \mathrm{Ad}_{\mathrm{L(H)}} y\) . Then \(\mathrm{Ad}(x^{-1}y)\) is the identity on L(H), hence \(x^{-1}y \in Z_{\overline{\mathrm{H}}}(\mathrm{H})\) and \(x \in Z_{\overline{\mathrm{H}}}(\mathrm{H})\) . H. By continuity, \(Z_{\overline{\mathrm{H}}}(\mathrm{H})\) is equal to C. The group C is closed in \(\overline{\mathrm{H}}\) and hence in G; therefore \(\overline{\mathrm{C}} \cap \overline{\mathrm{H}} \subset \mathrm{C}\) . Let \(x \in \mathrm{C}\) . There exists a sequence \((x_n)\) in H such that \(x_n\) tends to \(x\) . Then \(\mathrm{Ad}_{\mathrm{L(H)}} x_n\) tends to 1 and hence there exists a sequence \((y_n)\) in H such that \(y_n\) tends to \(e\) and \(\mathrm{Ad}_{\mathrm{L(H)}} y_n = \mathrm{Ad}_{\mathrm{L(H)}} x_n\) . Then \(y_n^{-1}x_n \in \mathrm{C} \cap \mathrm{H}\) and \(y_n^{-1}x_n\) tends to \(x\) .)
\item
  \(\mathbf{H} = \overline{\mathbf{H}}\) if and only if the centre of \(\mathbf{H}\) is closed in \(\mathbf{G}\) . (Use (b).)
\end{enumerate}

\begin{enumerate}
\def\labelenumi{\arabic{enumi}.}
\setcounter{enumi}{39}
\tightlist
\item
  Let \(\mathbf{H}\) be a finite-dimensional real Lie group and \(\mathbf{H}_0\) its identity component.
\end{enumerate}

\begin{enumerate}
\def\labelenumi{(\alph{enumi})}
\item
  Suppose that \(H_{0}\) satisfies condition (F) of Exercise 38, that \(H/H_{0}\) is finite and that the centre of \(H_{0}\) is compact. Let G be a finite-dimensional real Lie group and \(f:H\to G\) a continuous homomorphism with discrete kernel. Then \(f(\mathrm{H})\) is closed in G. (Reduce it to the case where H is connected. As the kernel of f is discrete, the centre of \(f(\mathrm{H})\) is the image under f of the centre of H (Lemma 1) and hence is closed in G. Apply Exercise 39 (c).)
\item
  Suppose that \(H_{0}\) is semi-simple and that \(H/H_{0}\) is finite. Then the image of H under a finite-dimensional continuous linear representation is closed in the linear group. (Apply (a) and Exercise 7 (b).)
\end{enumerate}

\begin{enumerate}
\def\labelenumi{\arabic{enumi}.}
\setcounter{enumi}{40}
\tightlist
\item
  Let \(G\) be a finite-dimensional connected real Lie group and \(\rho: G \to \mathbf{GL}(n, \mathbf{C})\) a continuous homomorphism with finite kernel. Then \(\rho((G, G))\) is closed in \(\mathbf{GL}(n, \mathbf{C})\) . (Using Exercises 7 (b) and 40 (a), reduce it to the case where \(G\) is the semi-direct product of a semi-simple group \(S\) and its radical \(R\) . Every element of \(\rho((G, R))\) is unipotent and hence \(\rho((G, R))\) is closed. The vector subspace \(V_1\) of fixed points of \(\rho((G, R))\) is \(\neq \{0\}\) . It is stable under \(\rho(G)\) . By considering \(\mathbf{C}^n / V_1\) arguing by induction on \(n\) and using the complete reducibility of \(\rho(S)\) , it is seen that it is possible to write
\end{enumerate}

\[
\mathbf {C} ^ {n} = \mathrm{V} _ {1} \oplus \dots \oplus \mathrm{V} _ {q}
\]

where each \(V_{i}\) is stable under \(\rho(S)\) and \(\rho((G,R))(V_{i})\subset V_{1}\oplus\cdots\oplus V_{i-1}\) for all i. On the other hand, \(\rho(S)\) is closed (Exercise 40 (b)). Finally, \(\rho((G,G))=\rho(S)\cdot\rho((G,R)).\)

\begin{enumerate}
\def\labelenumi{\arabic{enumi}.}
\setcounter{enumi}{41}
\tightlist
\item
  Let \(G\) be a finite-dimensional connected real Lie group. Suppose that
\end{enumerate}

G admits an injective continuous linear representation \(\rho: \mathrm{G} \to \mathbf{GL}(n, \mathbf{C})\) . Then G admits a linear representation \(\sigma: \mathrm{G} \to \mathbf{GL}(n', \mathbf{C})\) which is a homomorphism of G onto a closed subgroup of \(\mathbf{GL}(n', \mathbf{C})\) . (By Exercise 41, \(\rho((\mathrm{G}, \mathrm{G}))\) is closed in \(\mathbf{GL}(n, \mathbf{C})\) and hence (G, G) is closed in G. Let \(\rho: \mathrm{G} \to (\mathrm{G}, \mathrm{G})\) be the canonical morphism and \(\tau\) an injective linear representation with closed image \(\mathrm{G}/(\mathrm{G}, \mathrm{G})\) which is connected and commutative. Take \(\sigma = \rho \oplus (\tau \circ p)\) .)

